% TOP_PROBLEM: {"id": 2304026, "title": "Research Problems in Function Theory — Problem 4.26", "category_id": 8, "category": {"id": 8, "name": "analysis", "display_name": "Analysis", "description": "Limits, continuity, calculus, and function theory.", "slug": "analysis", "order_index": 8, "created_at": "2026-07-31T15:26:25.671Z"}, "status": "open", "problem_number": "AMR-022-4026", "set_id": 13}
% FIRST_POSTED: 2026-10-01
% ENTRY_KIND: statement_only
% UnsolvedMath ID 2304026; source catalogue code AMR-022-4026
% !TeX program = lualatex
% Definition and sources only; this file is not an LLM solution attempt.
\documentclass[11pt,a4paper]{article}
\usepackage[margin=25mm]{geometry}
\usepackage{fontspec}
\setmainfont{lmroman10-regular.otf}[BoldFont=lmroman10-bold.otf,
 ItalicFont=lmroman10-italic.otf,BoldItalicFont=lmroman10-bolditalic.otf]
\usepackage{amsmath,amssymb,mathtools}
\usepackage{unicode-math}
\setmathfont{latinmodern-math.otf}
\AtBeginDocument{\let\setminus\smallsetminus}
\usepackage{xurl,hyperref}
\hypersetup{colorlinks=true,urlcolor=blue}
\setcounter{secnumdepth}{0}
\setlength{\parindent}{0pt}
\setlength{\parskip}{5pt}
\setlength{\emergencystretch}{3em}
\begin{document}
\section*{Research Problems in Function Theory — Problem 4.26}\label{2304026}
{\small UnsolvedMath ID 2304026 · AMR-022-4026 · Analysis · AMR Open Problem Lists}
\subsection{Definitions and mathematical statement}
For an integer $n\geq0$, let $\mathcal P_n$ consist of the complex
polynomials
$p(z)=1+a_1z+\cdots+a_nz^n$ satisfying
$\operatorname{Re}p(z)>0$ whenever $|z|<1$. The degree may be
smaller than $n$. Positivity is required in the open disk; the
real part may vanish at boundary points.

Determine exactly
\[
 \Lambda_n=\max_{p\in\mathcal P_n}
       \int_0^{2\pi}|p(e^{i\theta})|^2\,d\theta
       =2\pi\max_{p\in\mathcal P_n}
                   \left(1+\sum_{j=1}^{n}|a_j|^2\right).
\]
The equality on the right is Parseval's identity and fixes the
normalization of the boundary integral. An extremal description
should identify the maximizing coefficient vectors, allowing for
rotations $p(z)\mapsto p(e^{it}z)$, which preserve both constraints
and the objective.

The analytic constraint can equivalently be expressed by the
nonnegativity, for all real $\theta$, of the trigonometric
polynomial
$1+\operatorname{Re}\sum_{j=1}^{n}a_je^{ij\theta}$.
Indeed, boundary nonnegativity gives interior positivity because
the value at the center is $1$. This formulation includes extremal
polynomials with boundary zeros of the real part. Coefficient
bounds for positive-real-part functions make the admissible
coefficient set bounded; its boundary-nonnegative formulation
makes it closed. Thus the maximum is an actual finite-dimensional
maximum, not merely a supremum. The target is its sharp value
and equality cases for every $n$.
\subsection{Short English statement}
Maximize the boundary square mean of a normalized polynomial with positive real part in the unit disk.
\subsection{Sources}
\begin{itemize}
\item[S1] ULAM.AI and the UnsolvedMath Contributors, supplied problems.json, record 2304026 (AMR-022-4026), AMR Open Problem Lists. Catalogue identity and source transcription; CC BY 4.0.
\url{https://huggingface.co/datasets/ulamai/UnsolvedMath}.
\item[S2] Hayman - Research Problems in Function Theory (2018), Problem 4.26. Original source indexed by AMR-022-4026.
\url{https://arxiv.org/abs/1809.07200}.
\end{itemize}
\end{document}
